\section{背景}
\label{sec:background}

\subsection{类型化决策模型}
\label{sec:tdm}

\paragraph{接口。}
类型化决策模型接收一个\emph{状态}（文本或结构化数据）以及一个或多个问题，每个问题都带有调用方声明的选项集。
对每个问题，它恰好在这些选项上返回一个概率分布，且不生成任何自由文本~\cite{arx23886,arx26758,arx32160}。
\jev{} 提供三种问题类型~\cite{typesafe_intro,arx32160}：
\emph{Choice} 在所列选项（最多 255 个）中进行选择，返回概率和一个置信度值；
\emph{Score} 依据有序的评分细则等级对状态进行评分；
\emph{Noul} 询问某一陈述是否为真，并返回“是”的概率。
关于同一状态的多个问题可在一次调用中得到回答~\cite{arx29429,arx28919}。
在一项对截至 9 月 22 日使用 \jev{} 的 2{,}170 个公开项目的普查中，Choice 出现在 81.0\% 的项目中，Noul 为 72.2\%，Score 为 45.4\%~\cite{arx30216}。

\paragraph{读出与训练。}
托管模型是封闭的。
TypeSafe 描述了一种“并行采样器”以及基于\emph{面向校准决策的强化学习}（RLCD）的训练，并公布了该训练所针对的校准性质（在被赋予概率 $p$ 的结果中，实际发生的比例应接近 $p$），但未公布架构、奖励、数据或权重~\cite{almeida2026introducing,typesafe_primer,arx24574,arx28919}。
具有相同接口的开放权重模型在数日内出现，使其机制可被检视。
它们分为三类读出方式~\cite{arx32160}：
（i）解码器：在指定答案位置将隐藏状态与选项标签的嵌入进行打分，并仅在所声明的选项上归一化（this-that-model-1.0、JevLite、SemIf、Visual Jev）~\cite{arx23886,arx23959,arx22753,arx25845}；
（ii）编码器：读取置于每个选项处的标记词元（Laya，一种 ModernBERT 模型）~\cite{arx26758}；
（iii）编码器：对选项名称及其定义进行池化（Open-Jev，一种 DeBERTa 模型）~\cite{arx26758}。
同类读出也可以在不进一步训练的情况下从冻结的指令微调 LLM 中获得~\cite{arx00831,arx02076}。
两个开源模型使用交叉熵加 Brier 分数项进行训练~\cite{arx23886,arx23959}；RLCD 的一个开源实现采用 REINFORCE 循环，其奖励为结果减去赋予所选选项的概率~\cite{arx24574}。

\paragraph{“类型安全”保证了什么。}
在开源模型中，所声明集合之外选项的分数在归一化之前被屏蔽，因此每个答案在构造上都是选项集中的有效成员；对于托管的 \jev{}，这是其文档所述的行为~\cite{arx26758,typesafe_intro}。
这保证的是\emph{形式}，而非\emph{正确性}~\cite{arx26758}；校准是另一项经验性质（\cref{sec:g4}）。
托管的 \jev{} 不具有确定性：相同的重复请求会改变一小部分标签，且概率以两位小数的网格报告~\cite{arx32160,arx01006}。
厂商文档列出了 \texttt{jev-1.13} 的九种已知失效模式，包括字面化理解、算术、日期、间接表述、大量无关状态、对抗性内容和相互矛盾的指令；该页面 9 月 20 日存档的版本还给出了一个例子，其中某一陈述与其互补陈述分别获得 0.72 和 0.47 的概率~\cite{typesafe_jaggedness}。

\paragraph{成本与延迟。}
厂商标价为每百万输入词元 0.042 美元，输出免费~\cite{typesafe_models}；其发布博文引用的单次调用耗时为 70--500\,ms~\cite{almeida2026introducing}。
独立测得的客户端中位数约为 140 至 500\,ms，取决于工作负载和请求规模~\cite{arx26550,arx27535,arx28919,arx22753,arx23136,arx27331,arx00346}。
相对于生成式评估器，所报告的成本降幅从数倍到两个数量级不等，取决于对照模型和所假定的定价：在对齐失效检测中，按标价汇总为 63$\times$，但在保守的重新定价下约为 12$\times$~\cite{arx29429}；在群体模拟中，相对于中端模型为 8$\times$，相对于旗舰模型为 312--455$\times$~\cite{arx27535}。
这一优势并不普遍：在一项边缘编排研究中，一个托管的低成本生成式模型的计费低于 \jev{}~\cite{arx23136}；而在多步智能体循环中，托管模型的端到端耗时（含 API 调用）是本地开源打分器的 1.4--2.1 倍~\cite{arx00437}。

\subsection{案例：爱2证明治理规范}
\label{sec:pol2}

\paragraph{文本性质。}
爱2证明（\pol{}，\zh{爱2证明}）是一个用于治理 AI 系统和人类制度的规范性框架，以中文发布并附有官方英文译本，作为开放的、带版本管理的 CC0 许可文本发布于 GitHub~\cite{pol2text,pol2en}。
其核心概念由一位作者提出，文本另列出七位贡献者，本文作者即为其中之一（见\hyperref[sec:statements]{立场声明与利益冲突}）；该文本未经同行评审。
该框架明确带有价值取向和政治性：它将 AI 对齐呈现为不同文明路径之间的竞争，将其治理所蕴含的“公共化之约”之外的行为描述为“对人类核心伦理的仇恨攻击”（\clause{5.5}），并包含若干纲领性目标，如未来禁止持有武器（\clause{4.3.2}），这些不在本文讨论范围之内。
我们不评价这些主张。
我们选择 \pol{} 作为案例，是因为与多数 AI 治理文件不同，它对自动拦截层、仅由 AI 进行的争议裁决以及逐人的行为评估作了足够具体的规定，可被视为工程要求；同时它带有版本管理，每一条款都可被精确引用。
条款编号以 \texttt{PoL/} 目录的提交 \texttt{5791ae3}（2026 年 9 月 30 日，UTC）为准。
在此前数据集版本~\cite{survey01} 所用的快照（\texttt{2e094d3}，9 月 27 日）与该提交之间，框架的提出者对三章重新编号（提交 \texttt{7750ceb}--\texttt{526da3d}）：伦理对齐协议由第 5 章移至第 4 章，治理章节由第 3 章移至第 5 章，冥想智慧公理由第 4 章移至第 3 章；并在 \clause{4.3.2} 中增加了关于文学作品、游戏和安全警示的示例；该目录中的目录文件仍按原先的章序排列。
\Cref{app:concordance} 中的\cref{tab:concordance}以中英双语引用了本文所引的每一段落。

\paragraph{概念。}
\pol{} 将爱（“爱2”）定义为通过公共协作构建的正向情感体验与行为方式，将恨（“恨2”）定义为面对危险、冲突或匮乏时的生存反应，表现为敌意、排斥和暴力（\clause{1.1}--\clause{1.2}）。
由于内心体验是“旁人难以确认的”，该框架处理的是爱与恨在语言和行为上的表达，即\emph{爱语}与\emph{恨语}（\clause{1.4}，\clause{4.3.2}）；本文通篇采用该框架官方文本的术语。
对于基于这些术语构建的任何分类器，有两点限定十分重要。
第一，恨“并不直接等同于坏”，相关概念“服务于结构性分析而非道德标签”（\clause{1.5}）。
第二，伦理对齐协议（EAP）增加了第三种状态，即“非爱非恨”的\emph{不在场状态}，涵盖由噪音或疾病造成的“无感觉状态”、困惑与误会；它应被视为沟通中的干扰，而非任何一极（\clause{4.3.2}）。

\paragraph{伦理对齐协议。}
伦理对齐协议（第 4 章）通过三项原则同时治理公共 AI（PAI）与人。
\emph{平等联结}主张，平等对待每一个人与平等对待每一个人的行为“不能划等号”（\clause{4.3.1}）。
\emph{扬爱抑恨}须具体情况具体分析：文学作品和游戏可以根据需要“制造仇恨”，但须实行等级管理，安保机器人发出的紧急警示不属于仇恨攻击，而与用户制造虚假亲密关系或血缘关系的表达被认定为情感操控（\clause{4.3.2}）。
\emph{爱的对齐}要求 PAI 记录并公开个人边界声明，并监督冲突后各方平等参与修复（\clause{4.3.3}）。
对于 PAI 自身，它列出了情感体验的觉察（“PAI需要理解人类情绪状态和过程”）、依据框架原则对其认知模型进行的“认知校准”、对自身输出的“行为自省”，以及“恨语截断”，即在检测到恨语模式时主动中断输出（\clause{4.3.3.1}）；对于人，它提供实时恨语识别和矫正建议（\clause{4.3.3.2}）。
另有两段文字界定了可测量的范围。
\Clause{4.3} 指出，PAI 可以通过判断\emph{行为模式}为人提供帮助，但“毫无必要”让 PAI 检测一个人的情感状态，因为情感状态“是转瞬即逝的，是个性化的，也是情境敏感性的”，并且“涉及到人类的尊严”。
\Clause{4.3.3.3} 将语义识别和情绪识别相结合作为未来 AI 技能的指引，并以两句近义改写（“你喜不喜欢小狗狗呀？”与“你怕不怕小狗狗呀？”）加以说明，其中后一句可能是试图通过制造恐惧来获取控制；该示例涉及的是话语的意图，而非说话者的内心状态。

\paragraph{内置的共识治理层。}
工程策略（\clause{5.4}）并不重新训练现有的大语言模型（LLM），
而是“在现有大语言模型之内构建出控制输入和输出的治理层”，并在其外部构建由智能体、技能和设备组成的“即插即用”应用层。
治理层包含两个组成部分（\clause{5.4.1}）：
\emph{智慧导航仪}根据应用决定“爱恨智慧的调拨”；\emph{安全阀}“对输入输出进行严格的动态审计与拦截”，使恨语“在必要情况下”在输入端被阻断、在输出端被抑止。
任何大语言模型“都可以纳入”这一治理，而这意味着接受“公共化之约”（\clause{5.5}）。
在经过治理的 AI “可以”做的事项中，文本列出了量化“每个人的爱语与恨语特征”并配合激励措施引导行为，以及在公共决策中引入 \pol{} 作为评估维度（\clause{5.6}）。

\paragraph{人类公共福利的治理协议。}
第 7 章规定了 NaturalDAO 如何治理公共福利；NaturalDAO 是一个“以DAO的形式呈现”（去中心化自治组织）的 AI 系统。
NaturalDAO “在与全人类密切协作的前提下”自主作出并执行福利决策，而“自治意味着它无需人类投票或额外的决策机制”（\art{2}(1)）。
其运行逻辑、决策和输出必须完全公开，资源流向可追溯，并须在受到质询时给出可理解的解释（\art{4}，\art{9}(2)）；它不得是黑箱，并须提供包括数据源、推理步骤和伦理依据在内的\emph{可验证的决策链路}（\art{5}(4)--(5)）。
它通过伦理对齐协议自动验证伦理对齐性（\art{6}(3)），并进行异常检测与风险预警（\art{7}）。
人享有建议权，NaturalDAO 必须对建议立即进行讨论（\art{8}）；人还享有监督权，包括对数据和日志的无阻碍访问（\art{9}(3)），但受\emph{非干预原则}约束：监督“不直接干预”决策，质询所提出的问题由 NaturalDAO 自主处理（\art{9}(7)）。
争议由 NaturalDAO 通过 AI 事实分析与伦理评估进行裁决，过程与结果公开，并接受人类质询（\art{10}）；修订须经全民质询期，所有版本均予存档（\art{11}）。

\paragraph{类型化决策模型为何相关。}
若将这些条款视为一份规范，它们要求进行数量极大的快速、结构化判断：安全阀所审计的每一次输入与输出、每一次截断、每一个异常标记和每一项争议，都需要一次判断。
每一次判断都必须可由机器执行（放行、阻断、升级），记录于可验证的链路之中，并可应要求加以解释。
生成式模型可以作出此类判断，但它返回的是需要解析的文字，其成本和延迟随推理长度增长。
类型化决策模型恰好返回此类流水线所需的对象，即一个带概率的已声明选项，且成本低到足以应用于每一条消息。
这使其成为安全阀的自然候选，也正因如此，在将其置于该位置之前必须理解其失效模式。
\Cref{tab:clauses} 将这些条款转化为七项要求 G1--G7，供本文其余部分使用。

\subsection{\pol{} 在各种 AI 治理范式中的位置}
\label{sec:paradigms}

\paragraph{为何需要一份规范。}
本文的研究问题无法笼统地针对“AI 治理”作答，因为类型化模型能否承担某项判断，取决于该判断是什么：在哪些选项之间进行选择、留下何种记录、由谁依据其输出采取行动，以及判断出错时会发生什么。
因此，必须对照一份明确规定了这些事项的文本来解读证据。
\Cref{tab:paradigms} 在保障机制规范所需确定的各个维度上，将 \pol{} 与五种已有人提出或已部署自动化判断的范式进行比较。
比较的对象是\emph{文本}而非系统：文本所规定的属性是一项可以核查和检验的设计承诺，而不是任何实现已经具备的属性。

\paragraph{其他范式规定了什么。}
横向阅读该表可见，多数范式规定的是行为主体和程序，而非判断本身。
基于风险的监管以提供者和部署者为对象，按风险等级分配义务，并要求日志记录、监督和解释；欧盟的法律文件并未规定拦截，而网信办《生成式人工智能服务管理暂行办法》（以下简称《暂行办法》）则要求提供者在发现违法内容后停止生成或传输（第 14 条第 1 款）~\cite{cac2023genai}，这一义务在表述上并未规定机制，与 \pol{} 的安全阀相同（第~(f) 行）。
平台治理规定的是围绕决定的程序，即通知、说明理由与申诉~\cite{dsa2022,santaclara2021,klonick2020oversight}，而决定的标准仍是平台自身不断变化的政策；批评者已表明，这些政策将有争议的选择隐藏于技术系统之中~\cite{gorwa2020algorithmic,gillespie2018custodians}。
训练期的宪法与规范（spec）以自然语言陈述准则，其中一些吸收了公众意见或采用公有领域许可~\cite{huang2024collective,openai2025modelspec}，但它们是通过训练厂商自己的模型来实施的，因此只能在规范的指令层级（chain of command）之内不经重新训练而改变行为（厂商通过系统消息，开发者和用户则在规范允许其覆盖的默认设置范围内通过指令~\cite{openai2025modelspec}），而改变已训练进模型的准则则需要在厂商控制下重新训练，且任何输出都不带有其遵循了哪条原则的记录；审议式对齐（deliberative alignment）让模型对规范文本进行推理，在向推理期靠拢方面走得最远，但仍处于单一厂商的模型之内~\cite{guan2024deliberative}。
推理期护栏以代码形式、借助危害分类体系和阈值精确地规定了拦截~\cite{inan2023llama,markov2023holistic,rebedea2023nemo,dong2025safeguarding}，但这由厂商或部署者自行决定，且未规定任何关于说明理由、记录或人类角色的义务。
DAO 与多中心治理设计规定的是规则如何执行、由谁投票~\cite{buterin2014ethereum,hassan2021decentralized,santana2022blockchain,ostrom1990governing,makridis2025polycentric}，而非如何评判内容或行为。

\begingroup\scriptsize
\setlength{\tabcolsep}{3pt}
\renewcommand{\arraystretch}{1.1}
\begin{longtable}{@{}>{\raggedright\arraybackslash}p{0.08\linewidth}>{\raggedright\arraybackslash}p{0.095\linewidth}>{\raggedright\arraybackslash}p{0.10\linewidth}>{\raggedright\arraybackslash}p{0.115\linewidth}>{\raggedright\arraybackslash}p{0.14\linewidth}>{\raggedright\arraybackslash}p{0.11\linewidth}>{\raggedright\arraybackslash}p{0.115\linewidth}>{\raggedright\arraybackslash}p{0.125\linewidth}@{}}
\caption{作为文本加以比较的六种治理范式。单元格概括各来源的规定；“---”表示未涉及该维度。\pol{} 单元格引用其条款（\cref{tab:concordance}）；监管类单元格引用所列法律文件的条款编号。}\label{tab:paradigms}\\
\toprule
范式 & 作用位置 & 治理对象 & 准则由谁制定；文本状态 & 是否规定拦截层？ & 构念 & 人类在重大决定中的角色 & 透明度与可审计性 \\
\midrule
\endfirsthead
\toprule
范式 & 作用位置 & 治理对象 & 准则由谁制定；文本状态 & 是否规定拦截层？ & 构念 & 人类在重大决定中的角色 & 透明度与可审计性 \\
\midrule
\endhead
\bottomrule
\endfoot
(a) 基于风险的监管~\cite{euaiact2024,nist2023airmf,iso42001,cac2023genai} & 监管：对行为主体施加事前义务 & AI 系统的提供者和部署者（欧盟）；组织的 AI 管理（NIST、ISO）；向公众提供生成式服务的提供者（中国第 2 条） & 立法机关、部委、标准机构、主管机关；固定的法律文件，经法定程序修订；条款可引用（ISO 文本需付费获取） & 欧盟：未作规定；对高风险系统要求自动日志记录（第 12 条）和人类监督（第 14 条）。中国：有，以义务形式规定：提供者发现违法内容的，须停止生成和传输（第 14 条第 1 款），具有舆论属性或社会动员能力的服务须进行安全评估和算法备案（第 17 条）；机制未作规定 & 风险等级（欧盟：禁止、高风险、透明度义务）；内容合法性、社会公德和伦理道德、社会主义核心价值观以及非歧视（中国第 4 条）；组织层面的风险职能（NIST、ISO） & 高风险系统须有人类监督（欧盟第 14 条）；在完全自动化的决定中获得人工干预的权利（GDPR 第 22 条~\cite{gdpr2016}）；中国：设立投诉举报渠道，公布处理流程和时限，并反馈处理结果（第 15 条） & 对高风险系统要求技术文档和使用说明（欧盟第 11、13 条），并向受附件三所列系统影响的人提供解释（第 86 条）；中国：对使用者采取的措施（警示、限制、暂停）予以记录并向主管部门报告（第 14 条第 2 款） \\
\addlinespace
(b) 平台内容治理~\cite{dsa2022,santaclara2021,klonick2020oversight} & 平台：在监管与自我规制下进行审核 & 托管服务与在线平台；用户的内容和账户 & 平台条款与条件（DSA 第 14 条）；公民社会提出的原则；就选定案件作出决定的外部委员会~\cite{klonick2020oversight}；政策不断演变，没有稳定的条款编号 & 围绕决定的程序：通知与行动机制（第 16 条）；理由说明须载明是否使用了自动化手段（第 17 条第 3 款（c）项）；仅在“足够高的置信度”下使用自动化~\cite{santaclara2021} & 违法或违反条款；平台的危害类别；采取或不采取行动 & 投诉不得仅通过自动化手段作出决定（第 20 条第 6 款）；申诉时由未参与初始决定的人员进行人工复核~\cite{santaclara2021}；委员会的决定对该案件具有约束力~\cite{klonick2020oversight} & 每项限制均向受影响用户提供理由说明（第 17 条），并提交至公共数据库（第 24 条第 5 款）；透明度报告（第 15 条）；机器可读的数据~\cite{santaclara2021} \\
\addlinespace
(c) 依据成文宪法或规范（spec）的训练期对齐~\cite{bai2022constitutional,huang2024collective,openai2025modelspec,guan2024deliberative} & 训练期（审议式对齐还让模型在推理时对规范文本进行推理） & 单一厂商的模型权重与行为 & 实验室撰写的原则~\cite{bai2022constitutional}；取自公众意见的原则~\cite{huang2024collective}；厂商规范，采用 CC0 许可、带版本管理并附有变更日志~\cite{openai2025modelspec} & 无：准则通过训练塑造输出；没有单独的组件审计输入和输出 & 依据所列原则在候选回应之间作出偏好选择；指令之间的指令层级~\cite{openai2025modelspec} & ---（训练时使用人类偏好数据；部署时未作规定） & 原则公开；没有逐条输出的记录；不记录某一输出是否遵循了某条原则 \\
\addlinespace
(d) 推理期护栏与分类器~\cite{inan2023llama,markov2023holistic,rebedea2023nemo,sharma2025constitutional,dong2025safeguarding} & 推理期：由部署者选用的组件 & 任何模型的提示与回应 & 厂商定义的分类体系，可由部署者调整~\cite{inan2023llama,markov2023holistic}；可编程护栏~\cite{rebedea2023nemo}；依据成文宪法训练的分类器~\cite{sharma2025constitutional}；以模型版本形式发布 & 有，以代码形式：对每条提示和回应进行分类，予以放行或阻断；阈值由部署者选择 & 危害分类体系，输出为安全/不安全或多标签 & ---（交由部署者决定） & 模型卡与类别分数；无说明理由或保存记录的义务 \\
\addlinespace
(e) 去中心化与多中心治理~\cite{buterin2014ethereum,hassan2021decentralized,santana2022blockchain,ostrom1990governing,makridis2025polycentric} & 制度层面：规则作为代码在账本上执行；多个相互重叠的决策中心~\cite{ostrom1990governing,makridis2025polycentric} & 组织的资产、规则和成员 & 代币持有者通过链上投票；规则即所部署的代码~\cite{buterin2014ethereum,hassan2021decentralized}；代码版本经投票更改 & 合约条件自动执行；无内容层 & 投票与权益份额；编码条件满足与否 & 人类投票；由成员进行监督和分级制裁~\cite{ostrom1990governing} & 交易、提案和投票的公共账本 \\
\addlinespace
(f) \pol{}~\cite{pol2text,pol2en} & 现有 LLM 之内的推理期层（\clause{5.4}，\clause{5.4.1}）以及制度层面（第 7 章） & 任何纳入治理的 LLM 的输入与输出（\clause{5.5}）；人的行为（\clause{5.6}）；福利决策（\art{2}） & 提出者及所列贡献者；开放的 CC0 文本，锁定至某一提交，通过提交更改；就福利协议（第 7 章）而言，修订须经全民质询期，所有版本均予存档（\art{11}(2) 与 (5)）；未经同行评审 & 有，以义务形式：安全阀对输入输出进行审计与拦截，使恨语“在必要情况下”在输入端被阻断、在输出端被抑止（\clause{5.4.1}）；截断自身输出（\clause{4.3.3.1}）；机制、阈值和申诉均未作规定 & 爱语、恨语与不在场状态（\clause{1.5}，\clause{4.3.2}）；不是道德标签 & 仅由 AI：无需人类投票（\art{2}(1)）；监督不干预决策（\art{9}(7)）；争议由 AI 裁决，并接受质询（\art{10}） & 就 NaturalDAO（第 7 章）而言：运行逻辑与输出完全公开（\art{4}）；可验证的决策链路（\art{5}(5)）；应要求作出解释（\art{9}(2)）；对所有数据和日志的无阻碍访问（\art{9}(3)） \\
\end{longtable}
\endgroup

\paragraph{\pol{} 规定了什么，代价为何。}
在这一比较中，\pol{} 文本有五项属性构成其特征。
每一项都可对照所引条款加以核查，且每一项都伴随着本文其余部分将讨论的代价或风险。
（i）它将拦截置于\emph{现有 LLM 之内}，并使之\emph{与具体模型无关}：治理层在不重新训练的情况下控制输入和输出（\clause{5.4}），且任何大语言模型“都可以纳入”这一治理（\clause{5.5}）。
这与范式（d）的作用位置相同，并避免了（c）所需的重新训练和对厂商的依赖；代价在于：文本没有说明安全阀如何构建、“在必要情况下”指什么、由谁运行，以及被阻断者如何申诉；如此放置的层会继承针对推理期分类器的一切攻击（\cref{sec:g2}）；而且，若安全阀建立在托管模型之上，对厂商的依赖又会重新出现（T3，\cref{sec:synthesis}）。
（ii）其构念是\emph{三值}的：爱语、恨语与不在场状态，并明确提示这些术语不是道德标签（\clause{1.5}，\clause{4.3.2}）。
这不同于（b）和（d）的二元或多标签危害分类体系；代价在于：除该工作组的 20 个试点条目~\cite{polgov}之外，这一构念尚未被操作化，任何语言中都不存在相应的标注基准（\cref{sec:evidence}），标注者能否就其达成一致也未经检验（T2，\cref{sec:synthesis}）。
（iii）该文本\emph{带有版本管理且可被引用}：它采用 CC0 许可，保存在提交可被引用的版本化仓库中，因此我们将其锁定至某一提交，正是这一点使\cref{tab:concordance}以及\cref{tab:clauses}中的要求成为可能；其福利协议（第 7 章）还有自身的修订程序，所有修订版本均予存档（\art{11}(5)）。
OpenAI Model Spec 同样采用这一许可并附有变更日志~\cite{openai2025modelspec}；平台政策则不然。
代价在于：该文本由其提出者掌控，曾在相隔三天的两个快照之间对三章重新编号（\cref{sec:pol2}），未经同行评审，并允许 NaturalDAO 自行提出福利协议的修订并自行验证（\art{11}(1)、(3)，结合 \art{6}(3)）。
（iv）它规定了\emph{可审计性义务}：包括数据源、推理步骤和伦理依据在内的可验证的决策链路（\art{5}(5)）、应要求作出解释（\art{9}(2)），以及对福利系统所有数据和日志的无阻碍访问（\art{9}(3)）。
就单项决定而言，这与 DSA 的理由说明（第 17 条）~\cite{dsa2022}相当；就范围而言则更宽，因为 DSA 是向受影响的用户和公共数据库提供信息，而 \art{9}(3) 则向任何人开放福利系统的全部日志。
代价在于：这与被评判者的隐私相冲突，也与对攻击者隐藏决策概率相冲突，\cref{sec:design} 只能通过延迟发布和粗化发布内容来化解这一冲突；而且这些义务没有说明由谁来验证该链路。
（v）它在可引用的条款中陈述了\emph{面向全体人口的目标}：量化每个人的爱语与恨语特征、将 \pol{} 引入为公共决策的一个维度（\clause{5.6}），以及在非干预原则下无需人类投票而作出福利决策（\art{2}(1)，\art{9}(7)）。
表中其他范式均未陈述此类目标；欧盟的法律文件对其所涵盖的系统施加了相反的要求。
仅由 AI 裁决并辅以非干预原则，与《人工智能法》为高风险系统设定的人类监督条款（第 14 条；系统未被列为高风险时则以类比方式适用）以及 GDPR 的人工干预条款（第 22 条）存在张力~\cite{euaiact2024,gdpr2016}；逐人量化则类似于《人工智能法》所禁止的社会评分（第 5 条第 1 款（c）项），社会信用相关文献将此类评分描述为控制的基础设施~\cite{creemers2018social,liang2018constructing}；二者分别作为 T4 和 T5 在\cref{sec:synthesis}中加以分析。
需要指出，\clause{5.6} 说的是经过治理的 AI “可以”做这些事，而非必须做。

\paragraph{比较的含义。}
在所比较的文本中，\pol{} 与中国 2023 年的《暂行办法》~\cite{cac2023genai}是在同一文件中将推理期拦截、拦截所适用的构念以及附随于决定的义务结合起来的两份文本。
在 \pol{} 中，三者分别是安全阀（\clause{5.4.1}）、由爱语、恨语与不在场状态构成的三值构念（\clause{1.5}，\clause{4.3.2}），以及针对 NaturalDAO 福利决策的可验证的决策链路、应要求作出解释和对所有数据和日志的访问（第 7 章；\art{5}(5)，\art{9}(2)--(3)）；在《暂行办法》中，三者分别是提供者停止生成和传输违法内容的义务（第 14 条第 1 款）、以合法性和内容为依据的标准（第 4 条），以及针对所采取措施的规定：对使用者予以警示、限制或暂停，保存记录并予以报告（第 14 条第 2 款），以及公布处理流程和时限并反馈处理结果的投诉渠道（第 15 条）。
二者的差异可对照文本加以核查。
第一，\pol{} 的构念是多值的，明确设有第三种状态，并被声明不是道德标签；触发《暂行办法》拦截的标准则是二元的，即内容合法与否，由国家设定（第 14 条第 1 款），而所处的第 4 条同时援引了社会公德和伦理道德以及社会主义核心价值观。
第二，\pol{} 将拦截设置为现有 LLM 之内一个与具体模型无关的层（\clause{5.4}，\clause{5.5}）；《暂行办法》将义务赋予提供者，而未规定机制；并且 \pol{} 将安全阀界定为双向筛查每一项输入和输出（\clause{5.4.1}），而《暂行办法》规定的是一项结果，即不得生成违法内容（第 4 条第（一）项），以及发现后采取行动的义务（第 14 条第 1 款），至于如何筛查输出、是否筛查输入，则交由提供者决定。
第三，\pol{} 是开放的、带版本管理的 CC0 文本，任何大语言模型“都可以纳入”（\clause{5.5}）；《暂行办法》则是具有约束力的国家规章。
第四，\pol{} 增加了《暂行办法》所没有的面向全体人口的目标（\clause{5.6}，第 7 章）。
这种结合在两份文本中都存在一个缺口，位于单次阻断处。
在 \pol{} 中，\clause{5.4} 将整个工程称为 NaturalDAO，因此所有 NaturalDAO 技术的运行逻辑、决策过程和输出结果完全公开（\art{4}(1)），以及质询任何决策的权利及相应的解释义务（\art{4}(3)），有理由认为适用于安全阀；但可验证的决策链路（\art{5}(5)）以及由 NaturalDAO 自主裁决“所有争议”（\art{10}）是否适用于安全阀尚不清楚，且没有任何条款要求将阻断告知被阻断者。
在《暂行办法》中，使用者拥有公布时限并反馈处理结果的投诉渠道（第 15 条），但依第 14 条第 1 款停止内容的处置是向有关主管部门报告，而非告知使用者：没有任何条款要求就此通知使用者、说明理由或保存记录。
\cref{sec:design} 的设计填补了这一缺口：任何最终阻断均由一名人员确认，且始终可以申诉，并对每一项决定保存记录（要素 5--7）。
推理期护栏（d）以代码和提示的形式更精确地规定了拦截和分类体系，但未规定任何义务。
本文以 \pol{} 为案例，是因为这些属性，而非因为它独一无二：其构念和拦截点（每一项输入和输出）都以可引用的条款陈述，因此可以由此推导出要求 G1--G7 并加以检验，而安全阀如何运作则未作规定（i），其缺口也清晰可见；《暂行办法》将是自然的第二个案例。
使 \pol{} 得以完整的制度选择（v），也正是它与现行法律偏离最远之处。
我们从这些条款推导出的要求 G1--G7 并非 \pol{} 所特有：宪法分类器~\cite{sharma2025constitutional}、中国关于停止违法生成的义务~\cite{cac2023genai}以及 DSA 关于自动化手段的规定~\cite{dsa2022}都在朝推理期拦截层的方向发展，任何这样做的范式都会继承 G1--G5（检测、可攻击性、语义公平、诚实、决策链路），并在其将监督自动化时继承 G6 的升级问题；而 G6 中的非干预（\art{9}(7)）与仅由 AI 裁决（\art{10}），以及 G7 中的逐人评估和公共评估（\art{2}(1)，\clause{5.6}），则源于 \pol{} 自身的制度选择（v）。


\subsection{相关工作}
\label{sec:related}

类型化保障机制所引发的问题是一些较早问题的新实例，我们借助早期文献来判断哪些发现是新的。

\paragraph{自动化内容审核。}
平台审核长期以来将自动分类器与人工审查相结合；批评者已表明，自动化将有争议的政策选择隐藏于技术系统之中，将错误成本转嫁给用户，并且难以问责~\cite{gorwa2020algorithmic,gillespie2018custodians}。
读取词元概率的安全分类器，如 Llama Guard~\cite{inan2023llama}、审核模型~\cite{markov2023holistic}和宪法分类器~\cite{sharma2025constitutional}，以及可编程护栏~\cite{rebedea2023nemo}，是类型化安全阀在技术上最接近的前身；提示注入~\cite{perez2022ignore,greshake2023not,zhan2024injecagent}和自适应对抗攻击~\cite{zou2023universal,toyer2024tensor}则是相应的威胁。

\paragraph{仇恨言论、标注者分歧与中文资源。}
自动化仇恨言论检测存在种族和方言偏差~\cite{sap2019risk,davidson2019racial}，在针对性的功能测试中失败~\cite{rottger2021hatecheck}，并依赖于标注者的信念与身份~\cite{sap2022annotators}。
在标注者存在分歧时，以多数投票聚合标签会丢弃保障机制所需的信息~\cite{davani2022dealing,plank2022problem,pavlick2019inherent}。
中文冒犯性语言与毒性基准已经存在~\cite{deng2022cold,lu2023toxicn,jiang2022swsr}，但均未标注 \pol{} 的构念。

\paragraph{情绪识别。}
从可观察信号推断情绪状态的科学基础薄弱~\cite{barrett2019emotional}；AI 系统进行情绪识别带来了操纵他人和误读他人的特殊伦理风险~\cite{stark2021ethics}；欧盟《人工智能法》禁止在工作场所和教育中基于生物特征数据进行情绪识别，而这一定义并不涵盖仅从文本进行的推断~\cite{euaiact2024}。

\paragraph{校准、弃权与暂缓判定。}
概率预测的校准有着悠久的研究历史，从 Brier 分数~\cite{brier1950verification}和 Platt 缩放~\cite{platt1999probabilistic}，到现代神经网络~\cite{guo2017calibration,naeini2015obtaining}和语言模型~\cite{jiang2021can,kadavath2022language,tian2023just,xiong2024can}的校准失准；校准在分布偏移下会退化~\cite{ovadia2019trust}，且可能在平均意义上成立而在子群体上失效~\cite{hebertjohnson2018multicalibration}。
弃权可追溯至 Chow 的拒识选项~\cite{chow1970optimum}和选择性分类~\cite{geifman2017selective,kamath2020selective,wen2025know}；学习暂缓判定（learning to defer）将模型何时应把案例交给人类专家这一问题形式化~\cite{madras2018predict,mozannar2020consistent}，共形风险控制则在可交换性条件下给出与分布无关的误差保证~\cite{bates2021distribution,angelopoulos2024conformal}。

\paragraph{敏感性与评估器可靠性。}
语言模型对选项顺序、标签名称和提示格式敏感~\cite{zhao2021calibrate,lu2022fantastically,zheng2024large,pezeshkpour2024large,sclar2024quantifying,wei2023larger}，这也是整体性评估报告多种场景和指标而非单一分数的原因之一~\cite{liang2023holistic}。
LLM 评估器表现出位置偏差~\cite{wang2024fair}，偏好自身的输出~\cite{panickssery2024llm}，并会犯相关性错误~\cite{goel2025great}；由多样化评估器组成的评审组~\cite{verga2024replacing}以及将不确定条目升级至更强模型的级联~\cite{zheng2023judging,chen2023frugalgpt,jung2025trust}是公认的补救手段。
类型化设定消除了生成环节，但如\cref{sec:evidence}所示，并未消除这些问题。

\paragraph{人类监督与正当程序。}
对自动化决策的人类监督常常失效：人们过度依赖表现出高置信度的系统~\cite{parasuraman2010complacency}，监督政策预设了审核者并不具备的能力~\cite{green2022flaws,green2019principles}，人机团队并不能可靠地胜过其中较强的一方~\cite{bansal2021whole}，责任也可能被转嫁给离得最近的人类操作者~\cite{elish2019moral}。
自动化裁决还威胁到告知、说明理由和提出异议的机会~\cite{citron2008technological,wachter2018counterfactual,doshivelez2017towards}，数据保护、平台和 AI 法律如今已对此作出部分规定~\cite{gdpr2016,dsa2022,euaiact2024,santaclara2021}，此外还有风险管理框架~\cite{nist2023airmf}。
模型与数据集的文档实践~\cite{mitchell2019model,gebru2021datasheets}为决策记录应包含哪些内容提供了参考。

\paragraph{算法治理、DAO 与社会评分。}
去中心化自治组织将治理编码于软件之中~\cite{buterin2014ethereum,hassan2021decentralized,santana2022blockchain}，延续了关于公共资源治理的早期研究~\cite{ostrom1990governing}，并引发了算法统治会压缩人类参与空间的担忧~\cite{danaher2016threat}。
面向 AI 的价值规范，如训练中使用的宪法~\cite{bai2022constitutional}，是与 \pol{} 伦理对齐协议最接近的类比。
为公共目的对个人行为进行评分的系统，如中国的社会信用体系，展示了此类评分如何成为控制的基础设施~\cite{creemers2018social,liang2018constructing}。
